%! Setting Up a Test for the Difference of Two Population Proportions — Unit 6, Lesson 6.10 — Homework
\documentclass[10pt]{article}
\usepackage{apstats-article}
\usepackage{apstats-boxes}
\newcommand{\TallMath}[1]{$\displaystyle #1\rule[-1.4em]{0pt}{3.2em}$}

\begin{document}

\pageheader{Unit 6, Lesson 6.10}{Homework}
\namedateperiod

\vspace{0.15in}

\begin{enumerate}

\item \textbf{College Club Membership.} A sociologist wants to test whether the proportion
of first-generation college students who join a campus club differs from the proportion of
non-first-generation students who do. Sample 1 (first-gen): 48 of 160 joined.
Sample 2 (non-first-gen): 72 of 200 joined.

\begin{enumerate}[label=\alph*.]
  \item State $H_0$ and $H_a$. \writelines{2}
  \item Compute $\hat p_1$, $\hat p_2$, and $\hat p_c$. \writelines{2}
  \item Verify all three conditions. \writelines{3}
\end{enumerate}

\item \textbf{Defect Rate.} A manufacturer claims that its new production line has a lower
defect rate than the old one. Old line: 36 defects in 400 units. New line: 20 defects in 400 units.

\begin{enumerate}[label=\alph*.]
  \item Define $p_O$ and $p_N$, then state $H_0$ and $H_a$ (one-sided). \writelines{2}
  \item Compute $\hat p_c$. \writelines{2}
  \item Verify conditions. \writelines{3}
\end{enumerate}

\end{enumerate}

\end{document}
